\section{Evidence Preservation: que los hechos puedan verificarse}

En un incidente grave, la memoria del equipo se altera por la presión, las conversaciones paralelas y los resultados posteriores. La preservación de evidencia no busca «preparar una versión para un litigio», sino que la investigación técnica, la protección de los usuarios, los seguros, la respuesta a reguladores y la revisión posterior partan del mismo conjunto de insumos verificables. La preservación debe comenzar lo antes posible y, al mismo tiempo, no debe obstaculizar la containment.

\subsection{Chain of custody, legal hold y timeline}

La \term{chain of custody} (cadena de custodia) registra quién recopiló, transfirió, consultó y almacenó la evidencia, cuándo y de qué manera, de modo que un revisor posterior pueda juzgar su integridad. Un hash ayuda a detectar cambios en un archivo, pero el hash por sí mismo no demuestra que el método de adquisición fuera correcto. Un \term{legal hold} (aviso de conservación legal) es el procedimiento por el cual el counsel ordena suspender la eliminación habitual y conservar la información que pueda ser pertinente; los ingenieros no deben declarar por su cuenta el privilege ni decidir el alcance de la conservación.

\lecturefigure{07-evidence-chain.png}{La calidad de la evidencia depende de una cadena continua de preserve, integrity, custody y timeline.}{Guía de respuesta a incidentes del NIST; subtítulos locales 21:00--25:40; redibujo conceptual.}

\begin{knowledgebox}{Lectura de la figura: la integridad proviene del proceso, no solo de las herramientas}
La etapa de preserve determina los logs, systems, snapshots y communications; integrity registra el hash, el timestamp y el acquisition method; custody restringe y audita el acceso, la transfer y el purpose; timeline ordena en el tiempo los hechos, las hipótesis, las decisiones y las correcciones posteriores. No basta con tomar una captura de pantalla de una conversación, ni debe permitirse que todos modifiquen directamente el mismo directorio de evidencia. Cada exportación debe tener un owner, un sistema de origen, una zona horaria y una política de retención.
\end{knowledgebox}

\begin{warningbox}{«Limpiar la escena» puede destruir a la vez la evidencia y la recuperación}
Eliminar cuentas, reinstalar sistemas o rotar todas las credenciales puede ser una containment necesaria, pero si no se registra ni se recopila nada, se pierde la root-cause evidence. A la inversa, retrasar la protección de los usuarios para hacer el análisis forense también puede ampliar el daño. El IC debe hacer que el technical lead y forensics/counsel definan con claridad qué acciones pueden ejecutarse de inmediato, cuáles requieren un snapshot y cuáles no pueden preservarse, pero deben registrarse con su justificación.
\end{warningbox}

\subsection{Decision log: separar hechos, análisis y autorización}

Un registro de calidad no consiste en conservar para siempre cada mensaje de Slack, sino en mantener un decision log estructurado: hechos técnicos actuales, incógnitas, análisis jurídico, opciones de negocio, decision owner, hora, fundamento y condiciones de nueva revisión. El \term{privilege} (privilegio de confidencialidad abogado-cliente), en el contexto de Estados Unidos, suele proteger comunicaciones confidenciales específicas realizadas para obtener o brindar asesoría jurídica; incluir a un abogado en un canal o marcar un documento como privileged no hace que todos los hechos de negocio desaparezcan ni que queden protegidos automáticamente.

\lecturefigure{09-decision-log.png}{Los hechos técnicos, el análisis jurídico, las decisiones de la alta dirección y el registro deben mantener sus límites dentro de una misma cadena.}{Subtítulos locales 21:00--25:40; guías del NIST y la SEC; redibujo conceptual.}

\begin{knowledgebox}{Lectura de la figura: registrar «quién decidió qué», no un consenso ficticio}
Los technical facts deben incluir evidence y confidence; el counsel analysis expone las obligaciones aplicables, los supuestos y el alcance de la confidencialidad; la executive decision enumera las options, la risk acceptance y el approver; record/audit conserva el owner, el rationale, la next review y la superseded decision. Cuando los hechos cambien más adelante, deben añadirse correcciones en lugar de sobrescribir las entradas anteriores. Así se evita que se suponga que el CSO tomó por sí solo una decisión jurídica, y también que los directivos afirmen después que nunca recibieron la información sobre el riesgo.
\end{knowledgebox}

\begin{importantbox}{Plantilla mínima de decision record}
\textbf{Time/owner}; \textbf{known facts and evidence links}; \textbf{unknowns}; \textbf{options}; \textbf{technical recommendation}; \textbf{legal/materiality analysis owner}; \textbf{business decision and approver}; \textbf{external communication}; \textbf{next checkpoint}. No escriba conjeturas en el campo de hechos, ni haga pasar el campo de recommendation por una decisión ya aprobada.
\end{importantbox}

\subsection{Resumen de la sección}

La cadena de evidencia responde «cómo lo sabemos», y el decision log responde «quién tomó qué decisión con base en qué hechos conocidos y desconocidos». Ambos protegen en conjunto la calidad de la investigación, los intereses de los usuarios y la responsabilidad de la organización; de ninguna manera sirven para fabricar memorandos que eximan de responsabilidad a una persona.
